% =====================================================================
% LAMPIRAN C: RINCIAN PEMROSESAN, EKSTRAKSI EMBEDDING, DAN FUSI DINI
% =====================================================================
\chapter{Rincian Pemrosesan, Ekstraksi \textit{Embedding}, dan Fusi Multimodal Dini}
\label{app:preprocessing-embedding}

Lampiran ini menampung rincian teknis tahap pemrosesan data multimodal, ekstraksi \textit{embedding} pada \textit{encoder} beku, dan skema fusi awal yang diringkas pada Bab~\ref{ch:metode}. Pemindahan ke lampiran dimaksudkan agar pembahasan Bab~\ref{ch:metode} tetap fokus pada keputusan metodologis tingkat tinggi, sementara langkah teknis tetap dapat diaudit secara mendetail.

\section*{C.1 Pemrosesan Data Gambar}
\label{app:img-preprocess}

Setiap gambar masukan diolah dengan urutan langkah berikut sebelum dilewatkan ke \textit{encoder} citra.

\begin{enumerate}
    \item \textbf{Pembacaan dan konversi}: berkas JPEG/PNG dibaca menggunakan \texttt{PIL.Image.open}, kemudian dikonversi ke \textit{mode} RGB.
    \item \textbf{Koreksi orientasi EXIF}: tag orientasi sensor kamera dibaca dan citra diputar ke orientasi normal agar foto vertikal tidak diumpankan dalam orientasi miring.
    \item \textbf{\textit{Resize}}: sisi pendek citra di-\textit{resize} ke 256 piksel (untuk DINOv2/DINOv3/Hiera) atau 512 piksel (untuk EVA-02) menggunakan interpolasi bilinear.
    \item \textbf{\textit{Center crop}}: pemotongan tengah ke ukuran target 224x224 atau 448x448.
    \item \textbf{Normalisasi per kanal}: piksel dinormalisasi dengan rerata $\mu = (0{,}485, 0{,}456, 0{,}406)$ dan simpangan baku $\sigma = (0{,}229, 0{,}224, 0{,}225)$ ImageNet sesuai Persamaan~\ref{eq:norm}.
    \item \textbf{\textit{Layout} tensor}: konversi ke tensor \texttt{float32} dengan \textit{layout} NCHW (\textit{batch, channel, height, width}).
\end{enumerate}

Karena \textit{encoder} citra dibekukan dan tidak dilatih ulang, augmentasi data tidak diterapkan pada saat ekstraksi \textit{embedding}.

\section*{C.2 Pemrosesan Data Teks}
\label{app:txt-preprocess}

Narasi laporan pada kolom \texttt{laporan} diolah dengan urutan langkah berikut.

\begin{enumerate}
    \item \textbf{Penghapusan URL}: \textit{regex} \texttt{https?://\textbackslash S+|www\textbackslash .\textbackslash S+} dipakai untuk menghapus tautan yang sering muncul pada laporan.
    \item \textbf{Normalisasi unicode}: konversi ke bentuk NFC.
    \item \textbf{Whitespace collapsing}: penggantian beberapa spasi atau \textit{tab} menjadi satu spasi.
    \item \textbf{Pembersihan karakter kontrol}: penghapusan karakter non-cetak.
    \item \textbf{Pemotongan}: pembatasan panjang teks ke 256 token, yang mencakup sekitar 98,5\% laporan tanpa pemotongan.
    \item \textbf{Pemberian prefiks}: untuk \textit{multilingual-E5-Large-Instruct}, teks diberi prefiks \texttt{Instruct: Classify Indonesian civic report\textbackslash nQuery: } sebelum tokenisasi, mengikuti konvensi varian \textit{instruct}.
    \item \textbf{Tokenisasi}: setiap teks ditokenisasi menggunakan tokenizer bawaan masing-masing \textit{encoder} (SentencePiece atau WordPiece) dengan parameter \texttt{max\_length=256}, \texttt{padding="max\_length"}, dan \texttt{truncation=True}.
\end{enumerate}

Tidak diterapkan \textit{stopword removal} atau \textit{stemming}, karena \textit{Transformer encoder} sudah menangani morfologi kata secara natif.

\section*{C.3 Encoder Citra yang Dievaluasi}

Empat \textit{encoder} citra dipilih sebagai kandidat ablasi, semuanya menghasilkan vektor \textit{embedding} berdimensi 1024. Daftar \textit{encoder} dirangkum pada Tabel~\ref{tab:image-encoders}.

\begin{table}[htbp]
\centering
\caption{Daftar \textit{encoder} citra yang dievaluasi.}
\label{tab:image-encoders}
\small
\begin{tabular}{p{2.5cm} p{5.0cm} c c c}
\toprule
\textbf{Model} & \textbf{Model ID} & \textbf{Dim} & \textbf{Input} & \textbf{Framework} \\
\midrule
DINOv2-L & dinov2-large & 1024 & 224 & PyTorch \\
DINOv3-ViT-L & dinov3-vitl16 & 1024 & 224 & PyTorch \\
Hiera-L & hiera\_large\_224 & 1024 & 224 & PyTorch \\
EVA-02-L & eva02\_large\_patch14\_448 & 1024 & 448 & PyTorch \\
\bottomrule
\end{tabular}
\end{table}

Vektor representasi diambil dari token CLS atau rata-rata token \textit{patch} sesuai konvensi masing-masing \textit{encoder}, kemudian dinormalisasi L2.

\section*{C.4 Encoder Teks yang Dievaluasi}

Empat \textit{encoder} teks dipilih sebagai kandidat ablasi, mencakup \textit{encoder} multibahasa generik dan \textit{encoder} khusus Bahasa Indonesia. Daftar dirangkum pada Tabel~\ref{tab:text-encoders}.

\begin{table}[htbp]
\centering
\caption{Daftar \textit{encoder} teks yang dievaluasi.}
\label{tab:text-encoders}
\small
\begin{tabular}{p{3.2cm} p{5.0cm} c p{2.8cm}}
\toprule
\textbf{Model} & \textbf{Model ID} & \textbf{Dim} & \textbf{Pool Strategy} \\
\midrule
mE5-Large-Instruct & multilingual-e5-large-instruct & 1024 & Mean masked \\
BGE-M3 & bge-m3 & 1024 & CLS \\
IndoBERT-Large-p2 & indobert-large-p2 & 1024 & CLS \\
Cendol-mT5-Large & cendol-mt5-large-inst & 1024 & Mean (encoder) \\
\bottomrule
\end{tabular}
\end{table}

\textit{Output} setiap \textit{encoder} dinormalisasi L2 setelah \textit{pooling}.

\section*{C.5 Pipeline Inferensi Batch dan Caching Embedding}

Inferensi \textit{batch} dilakukan dengan \textit{loop} di atas \texttt{DataLoader} PyTorch yang membaca masukan secara \textit{streaming} dari direktori gambar atau \texttt{Dataset} teks. Konteks \texttt{torch.no\_grad()} dan \texttt{torch.autocast(dtype=torch.float16)} dipakai untuk efisiensi memori. Vektor keluaran dikonversi ke \texttt{numpy} \texttt{float32} sebelum dituliskan ke berkas \textit{cache}.

\textit{Embedding} hasil ekstraksi disimpan pada berkas terkompresi yang dikelompokkan per \textit{encoder} dan per partisi data, dengan identitas baris yang sejajar dengan label. Strategi \textit{caching} ini memungkinkan eksperimen ablasi pada algoritma klasifikasi dilakukan tanpa perlu menjalankan ulang \textit{encoder}, sehingga setiap kombinasi (\textit{encoder} citra, \textit{encoder} teks) dapat dievaluasi dengan biaya komputasi yang sangat rendah.

\section*{C.6 Rincian Fusi Multimodal Dini}
\label{app:fusion-detail}

Penelitian ini mengadopsi skema fusi awal (\textit{early fusion}) berbasis konkatenasi, sebagaimana metodologi yang dipublikasikan oleh \citet{hanifPenugasanOtomatis2024}. Sebelum konkatenasi, setiap \textit{embedding} dinormalisasi L2 secara independen agar magnitudonya seragam dan tidak ada modalitas yang mendominasi karena perbedaan skala absolut. Normalisasi L2 didefinisikan sebagai $\mathbf{x}_{\text{norm}} = \mathbf{x} / \|\mathbf{x}\|_2$.

Vektor fusi akhir adalah konkatenasi langsung sesuai Persamaan~\ref{eq:fusion-method}:

\begin{equation}
    \mathbf{e}_{\text{fuse}} = [\mathbf{e}_{\text{img}}^{\text{L2}};\, \mathbf{e}_{\text{txt}}^{\text{L2}}] \in \mathbb{R}^{2048}.
    \label{eq:fusion-method}
\end{equation}

Vektor inilah yang menjadi masukan tunggal bagi algoritma klasifikasi CatBoost.

Sebagaimana dibahas pada Bagian~\ref{sec:pca}, PCA dapat dipakai untuk mereduksi dimensi \textit{embedding}. Pada penelitian ini PCA \textbf{tidak diterapkan} dengan dua pertimbangan, yaitu (1) CatBoost berbasis pohon keputusan \textit{oblivious} sehingga mampu menangani fitur padat berdimensi tinggi tanpa penalti \textit{curse of dimensionality} yang berarti, dan (2) reduksi dimensi linear seperti PCA dapat menghilangkan informasi diskriminatif yang justru diperlukan untuk membedakan kategori dinas yang dekat secara semantik.
